\scanpage{72}
% archive sheet number: 71

$\underline{Q}_p$ faisceaux constructibles sur $k$ vers la catégorie $\mathrm{Fcriso}(k)$. Pour un objet \deleted{$E$}$(E,F_E)$ de cette dernière, les conditions suivantes sont équivalentes\inserted{\deleted{\unsure{k|x}}\handwritten{(}si $k$ parfait\deleted{x}\handwritten{)}}: a) $(E,F_E)$ appartient à l'image essentielle du foncteur précédent; b) il existe un réseau $M\subset E$ stable par $F_E$ \deleted{sur} \uline{\unsure{\deleted{l}e|\deleted{e}l}t} par $F_E^{-1}$ (i.e. tel que $F_E$ induise un \uline{automorphisme} de $M$); c) la famille des i\unsure{\handwritten{t}|\deleted{e}\handwritten{t}|\deleted{o}\handwritten{t}|\deleted{a}\handwritten{t}}érés $F^{\unsure{n|m}}$, $n\quad\underline{Z}$, est bornée ; d) $E$ et $\overset{\handwritten{\vee}}{E}$ sont effectifs.

Les F-isocristaux en question sont appelés ``trivaux'' ou ``\uline{de pente nulle}''. On trouve ainsi un \uline{épimorphisme de gerbes}
\diagram{\begin{tikzcd}\underline{G}(k) \arrow[r,thick] & \handwritten{\underline{\pi}_{p}}(k)\ ,\end{tikzcd}}
\margin{\handwritten{$\Pi$\par majuscule?}}
où $\handwritten{\underline{\unsure{\Pi|\pi}}}_p(k)$ est la gerbe sur $\underline{Q}_p$ \deleted{des représentations continues} associée à la catégorie tannakienne des $\underline{Q}_p$-faisceaux constructibles, équivalente \deleted{à l'enveloppe alg} (une fois choisie une clôture algébrique $\bar{k}$ de $k$) à la gerbe triviale définie par l'enveloppe algébrique \deleted{du} $p$-adique du groupe pro-fini $\handwritten{\pi}=\operatorname{Gal}(\handwritten{\bar{k}}/k)$. Tenant compte de l'homomorphisme d'augmentation de \unsure{\deleted{G}T|\deleted{T}G}ate, on trouve même un homomorphisme
\diagram{\begin{tikzcd}\underline{G}(k) \arrow[r,thick] & \handwritten{\underline{\pi}}_p(k)\ \mathrm{x}\ \handwritten{[}\underline{G}_m\handwritten{]}_{\underline{Q}_p}\end{tikzcd}}
où $\handwritten{[}\underline{G}_m\handwritten{]}$ est la gerbe triviale définie par $\underline{G}_m$ . On vérifie facilement que c'est encore un épimorphisme (i.e. que le foncteur
\diagram{\begin{tikzcd}(V_i)_{\handwritten{i}} \arrow[r,maps to,thick] & \handwritten{\sum}V_i\otimes_{\underline{\unsure{\deleted{Z}Q|\deleted{Q}Z|\not Q}}_p}\underline{Q}_p(i)\end{tikzcd}}
allant de la catégorie des \deleted{\unsure{a|e}b} $\underline{Q}_p$-faisceaux constructibles \inserted{\handwritten{gradu\unsure{é|i}s}} sur $k$ vers $\mathrm{Fcriso}(k)$ est pleinement fidèle). Les F-isocristaux de la forme précédente pourraient être appelés les F-isocristaux \uline{de niveau zéro} .

\handwritten{6.} Considérons maintenant un homomor\unsure{\^{p}|\deleted{ô}p|\deleted{p}ô}\unsure{\deleted{i}h|\.{h}|\deleted{h}i}isme de corps\textsubscript{x}
\diagram{\begin{tikzcd}k \arrow[r,thick] & k\handwritten{\prime}\ ,\end{tikzcd}}
d'où un homomorphisme de catégories tannakiennes sur $\underline{Q}_p$
\diagram{\begin{tikzcd}\mathrm{Fcriso}(k) \arrow[r,harpoon,thick] & \mathrm{Fcriso}(\deleted{\unsure{i|t|k}}\handwritten{k}\prime)\ ,\end{tikzcd}}
et un homomorphisme de gerbes tannakiennes sur $\underline{Q}_p$ en sens inverse


